% Appendix to Section 6: proofs and details for many schemas without labels.
% Sources: research/tracks/untagged/notes-final.md (with referee.md and its verification log),
% research/tracks/single/notes-final.md section 10 (Lemma H.1, Theorems H.2--H.5).
\section{Proofs and details for Section~\ref{sec:many}}\label{app:many}

Notation is that of \cref{sec:many:setting}. ``Template'' means a template in $\DT$; ``covering $X$'' means having every
element of $X$ as an instance. The finitariness facts of \cref{prop:many:W} are used throughout.

%=====================================================================================================================
\subsection{The setting}\label{app:many:setting}

\begin{proof}[Proof of \cref{lem:many:cautious}]
(a) $\manyRstar=\bigcup_i\inst(\sigma_i)\in\Hk{k}$ because $k\ge k'$; it contains $D$ and avoids $N$. Larger $D$ or $N$
shrink the version space, so they enlarge its intersection (\cref{thm:setting:vs}(a), \cref{lem:setting:onesided}(b)).
(b) Each displayed union is in $\Hk{k}$ (at most $k$ blocks), contains $D$ and avoids $N$, so it lies in
$\manyVS_k(D,N)$; hence $\Acc_k(D,N)$ is contained in the displayed intersection. Conversely let
$h=\bigcup_l\inst(\tau_l)\in\manyVS_k(D,N)$. Assign each datum to the first slot covering it; this partitions $D$ into at
most $k$ nonempty blocks $B\subseteq\inst(\tau_l)$. By \cref{prop:many:W}(b), $\tau_l\gen T_B$ for some $T_B\in\Min(B)$, and
$\inst(T_B)\subseteq\inst(\tau_l)\subseteq h$ avoids $N$, so $T_B\in\Min_N(B)$ and $h\supseteq\bigcup_B\inst(T_B)$. So every
member of the version space contains a displayed union, and the displayed intersection is contained in $\Acc_k(D,N)$.
Decidability: there are finitely many partitions, each $\Min(B)$ is finite and computable, and avoidance of a finite $N$
and membership of a query are decided by matching (\cref{thm:setting:matching}).
\end{proof}

With $N=\emptyset$ and $\Acc(B):=\bigcap\manyCov(B)=\bigcap_{T\in\Min(B)}\inst(T)$ (and $\Acc(\emptyset):=\emptyset$), part (b)
says, since the choices $T_B$ are independent across blocks:
\begin{equation}\label{eq:many:partition}
q\in\Acc_k(D)\iff\forall\pi\ \exists B\in\pi:\ q\in\Acc(B),
\end{equation}
where $\pi$ ranges over the partitions of $D$ into at most $k$ nonempty blocks.

\begin{lemma}[monotonicity of coherence]\label{lem:many:mono}
\status{proved}\src{untagged Lemma 5.1}
$X$ is $d$-coherent iff some $T\in\manyCov(X)$ is not $d$-refuted. $d$-coherence is downward closed ($\emptyset\neq
Y\subseteq X$ and $X$ $d$-coherent imply $Y$ $d$-coherent) and antitone in $d$.
\end{lemma}
\begin{proof}
If $T\in\manyCov(X)$ is not $d$-refuted, then $T\gen T_0\in\Min(X)$, and $\inst(T_0)\subseteq\inst(T)$ is not $d$-refuted
either; the converse holds as $\Min(X)\subseteq\manyCov(X)$. A covering template of $X$ covers $Y$. And
$\Ref_d\subseteq\Ref_{d+1}$.
\end{proof}

%=====================================================================================================================
\subsection{Positive data: bounds, spare slots and head classes}\label{app:many:bound}

\begin{proof}[Proof of \cref{prop:many:bound}]
$D$ is a finite union of ground templates and avoids $N$. For the second claim take a true non-ground $\sigma$: the class
contains $\inst(\sigma)$ and all its finite subsets, all inside $\manyTr$, on which the oracle answers identically. So the
learner is a function of the text, and Gold's locking-sequence argument for superfinite classes applies
\citep{gold1967language}; it needs no computability.
\end{proof}

\begin{proof}[Proof of \cref{thm:many:splits}]
(a) $h\subseteq\manyRstar$ and $N\cap\manyRstar=\emptyset$ give $h\cap N=\emptyset$, and $D\subseteq h\in\Hk{k}$.
(b) Let $h=\bigcup_l\inst(\tau_l)\in\manyVS_k(D,N)$ and $L_j:=\{l:\tau_l\text{ covers a datum of }D_j\}$. Each $L_j$ is
nonempty because $D_j$ is. By cross-separation a slot covering data of two targets would meet $N$, so the $L_j$ are
pairwise disjoint and $|L_i|\le k-(k'-1)=m$. Every datum of $D_i$ is covered by a slot in $L_i$, and these slots avoid
$N$, so by hypothesis $\bigcup_{l\in L_i}\inst(\tau_l)\supseteq\inst(\sigma_i)$; hence $h\supseteq\inst(\sigma_i)$.
For the second claim, $h:=\bigcup\inst(F)\cup\bigcup_{l\neq i}\inst(\sigma_l)$ uses at most $m+k'-1=k$ slots, contains $D$,
avoids $N$ (the $\sigma_l$ lie in $\manyRstar$) and misses $q$. If the instance sets are pairwise disjoint and the first
condition fails, some admissible family misses some $q\in\inst(\sigma_i)$, which lies in no other instance set, so the
second claim applies: the condition is necessary and sufficient.
\end{proof}

In general a target's slots may be shared with others, provided every shared slot contains the whole target.

\begin{theorem}[slot accounting]\label{thm:many:slots}
\status{proved}\src{untagged Thm 4.1}
Fix a target $i$, a bound $k$ and a sound $N$, and let $\bar D_i:=D\setminus D_i$. Assume (Abs): every $N$-avoiding
template covering a datum of $D_i$ and one of $\bar D_i$ contains $\inst(\sigma_i)$. Let $c_i$ be the least number of
$N$-avoiding templates not containing $\inst(\sigma_i)$ that cover $\bar D_i$. (a) If every family of at most $k-c_i$
$N$-avoiding templates covering $D_i$ has union $\supseteq\inst(\sigma_i)$, then $\inst(\sigma_i)\subseteq\Acc_k(D,N)$.
(b) If some such family and some optimal cover of $\bar D_i$ both miss $q\in\inst(\sigma_i)$, then $q\notin\Acc_k(D,N)$.
\end{theorem}

\begin{corollary}[slot accounting for PA]\label{cor:many:paslots}
\status{proved; computed}\src{untagged Cor 4.2}
Let the targets be $\Q$'s axioms Q1--Q7 (numbered as in \cref{sec:setting:examples}) as ground sentences in
closure-normal form ($\neg Sp=0$; $Sp=Sp'\to p=p'$; $\neg p=0\to\exists y\,p=Sy$; $p+0=p$; $p+Sp'=S(p+p')$; $p\cdot0=0$;
$p\cdot Sp'=p\cdot p'+p$) and $T_{\Ind}$, so $k'=8$. (Abs) holds for induction for every motive: a template covering a
$\Q$-axiom and an induction instance lies above $F_0$, or above $F_0\to F_1$ for Q2 and Q3 ($F_0,F_1$ $0$-ary formula
metavariables), and both contain $\inst(T_{\Ind})$. On positive data $c=4$ (Q4--Q7 share $z_0=z_1$; Q1, Q2, Q3 stay
alone), so at $k=8$ induction gets four slots: an induction instance whose motive has $x$ free and a main
connective unseen in the data is accepted iff the induction data cannot be split into four blocks each root-homogeneous
or all-vacuous (without vacuous motives: iff at least five main connectives occur). With the three negatives $0=S0$, $(p+0)=0$, $(p\cdot0)=S0$,
$c=7$, induction gets one slot, and the tagged anchor condition \evR{}+\evN{} suffices. (In the Sub-encoded presentation
of induction of the earlier analysis of untagged PA, $c=1$ \src{prior pa-untagged}.)
\end{corollary}
The computed check of $c=4$ and $c=7$ (24/24 cases)\src{untagged u2} holds by construction given (Abs) and the anchor theorem
for $T_{\Ind}$ (\cref{thm:zf:indanchor}), so it is a consistency check of the implementation, not independent evidence.

\begin{proof}[Proof of \cref{thm:many:slots}]
(a) Let $h\in\manyVS_k(D,N)$ with slots $\tau_l$. If some slot contains $\inst(\sigma_i)$ we are done. Otherwise, by (Abs),
no slot covers both a datum of $D_i$ and a datum of $\bar D_i$. The slots covering $\bar D_i$-data form an $N$-avoiding
cover of $\bar D_i$ by templates not containing $\inst(\sigma_i)$, so there are at least $c_i$ of them, and they are
disjoint from the slots covering $D_i$. So at most $k-c_i$ slots cover $D_i$, and by hypothesis their union contains
$\inst(\sigma_i)$. (If $k\le c_i$, this shows that some slot of every $h$ contains $\inst(\sigma_i)$.)
(b) $\bigcup\inst(F\cup G)$ has at most $k$ slots, contains $D$, avoids $N$ and misses $q$.
\end{proof}

\begin{proof}[Proof of \cref{cor:many:paslots}]
\emph{(Abs).} At the top level of a sentence no bound variable is in scope, so a metavariable whose pattern occurrence
is not under a binder is 0-ary. A $\Q$-axiom with root $\neg$ or $=$ and an induction instance (root $\to$) have no common
root, so their minimal covering template is the bare $F_0$. Q2 and Q3 have root $\to$; their antecedents ($=$, resp.\
$\neg$) differ in head from the induction antecedent ($\wedge$), and their consequents ($=$, resp.\ $\exists$) from the
induction consequent ($\forall$); the two slots are closed and not equal in either datum, so the minimal covering template
is $F_0\to F_1$. This holds for every motive. Every covering template lies above one of these (\cref{prop:many:W}), and
both contain $\inst(T_{\Ind})$.
\emph{$c=4$ with $N=\emptyset$.} Any template covering two of Q1, Q2, Q3, or one of them and one of Q4--Q7, lies above
$F_0$ or $F_0\to F_1$ (compare the roots and, for Q2 and Q3, the heads of antecedents and consequents), hence contains
$\inst(T_{\Ind})$. So Q1, Q2 and Q3 need a slot each and Q4--Q7 at least one; $z_0=z_1$ covers Q4--Q7 and contains no
induction instance. \emph{$c=7$ with $N=\{0=S0,(p+0)=0,(p\cdot0)=S0\}$.} Q4 and Q5 have the lgg $p+z_0=z_1$ (common prefix
$=,+,p$; closed slots), which has the instance $(p+0)=0\in N$; Q6 and Q7 have $p\cdot z_0=z_1\ni(p\cdot0)=S0$; the other
pairs among Q4--Q7 have $z_0=z_1\ni0=S0$. All covering templates of these pairs lie above the lggs, so Q4--Q7 need a slot
each. ($F_0$ is now refuted, $F_0\to F_1$ is not, so (Abs) still holds.) Then induction gets $k-7=1$ slot, and by
\cref{thm:many:slots} it is learned exactly iff its data contain an anchor, which for $T_{\Ind}$ is \evR{}+\evN{}
(\cref{thm:zf:indanchor}).
\emph{The four-slot criterion.} Let $q$ be an induction instance whose motive has $x$ free and has a main connective that
occurs in no datum. By \cref{thm:zf:indanchor}, a block of induction data is a non-anchor iff its motives all have the same
root $f$ or are all vacuous; then $T_{\Ind}[P:=\lambda x.f(\dots)]$, resp.\ $T_{\Ind}[P:=\lambda x.A]$, covers the block and
misses $q$. Let $\chi$ be the least number of blocks in a partition of the induction data into blocks that are
root-homogeneous or all-vacuous (without vacuous motives: the number of main connectives occurring). If $\chi\le4$, these
templates for an optimal partition, together with the cover $\{$Q1, Q2, Q3, $z_0=z_1\}$ of the $\Q$-data, form a hypothesis
with $8$ slots missing $q$ (\cref{thm:many:slots}(b)). If $\chi>4$, every family of at most four templates covering the
induction data has, by partition into first covering members, an anchor block, whose template contains $\inst(T_{\Ind})$
(\cref{thm:many:slots}(a)).
\end{proof}

\begin{corollary}[thresholds in terms of failure families]\label{cor:many:failure}
\status{proved}\src{untagged Cor 3.3}
Call $T$ a \emph{proper specialization} of $\sigma$ if $\inst(T)\subsetneq\inst(\sigma)$. A \emph{failure family}
$\manyFail(\sigma)$ is a set of proper specializations such that every finite non-anchor $X\subseteq\inst(\sigma)$ lies in
$\inst(F)$ for some $F\in\manyFail(\sigma)$. (Examples: for first-order patterns the head specializations $\sigma[x\mapsto
f(\bar z)]$ and the identifications $\sigma[y\mapsto x]$, by Plotkin's \evR{}+\evD{}; for $T_{\Ind}$ the nine root
specializations $T_{\Ind}[P:=\lambda x.f(\bar P(x))]$, $f\in\{=,<,\neg,\wedge,\vee,\to,\leftrightarrow,\forall,\exists\}$
(eight in the untagged prototype, whose arithmetic has no $<$), and the vacuous one $T_{\Ind}[P:=\lambda x.A]$.) Let $m=k-k'+1$.
(a) If $D_i\subseteq\bigcup_{j\le m}\inst(F_j)$ with $F_j\in\manyFail(\sigma_i)$, and some $q\in\inst(\sigma_i)\setminus
\bigcup_j\inst(F_j)$ lies in no other $\inst(\sigma_l)$, then $q\notin\Acc_k(D,N)$ for every sound $N$.
(b) Suppose $N$ cross-separates $D$, $N$ \emph{target-separates} $D_i$ (it meets every template that covers a datum of
$D_i$ and whose instance set neither contains nor is contained in $\inst(\sigma_i)$), and the instance sets are pairwise
disjoint. Then $\inst(\sigma_i)\subseteq\Acc_k(D,N)$ iff $D_i$ is not covered by $m$ proper specializations of $\sigma_i$
whose union misses an instance of $\sigma_i$. If moreover no $m$ members of $\manyFail(\sigma_i)$ cover $\inst(\sigma_i)$,
this holds iff $D_i$ is not covered by $m$ members of $\manyFail(\sigma_i)$.
\end{corollary}
\begin{proof}
(a) Apply \cref{thm:many:splits}(a) to $h=\bigcup_j\inst(F_j)\cup\bigcup_{l\neq i}\inst(\sigma_l)$, which has $k$ slots and
satisfies $D\subseteq h\subseteq\manyRstar$.
(b) If exactness for $i$ fails, then by \cref{thm:many:splits}(b) (which uses cross-separation) some $N$-avoiding family
of at most $m$ templates covers $D_i$ and misses some $q\in\inst(\sigma_i)$. No member contains $\inst(\sigma_i)$, so by
target-separation each member covering a $D_i$-datum is a proper specialization. Each block $X_j=D_i\cap\inst(T_j)$ is
covered by a template not containing $\inst(\sigma_i)$, so it is not an anchor and lies in some $\inst(F_j)$,
$F_j\in\manyFail(\sigma_i)$. Conversely, a cover by proper specializations that misses $q$ gives $q\notin\Acc_k(D,N)$ by
\cref{thm:many:splits}(b) and disjointness, and a cover by $m$ failure templates misses some $q$ under the extra hypothesis,
so (a) applies.
\end{proof}

On positive data ($N=\emptyset$) target $i$ is learned exactly if every partition of $D_i$ into at most $k$ blocks (at most
$k-c_i$ under (Abs)) has an anchor block, and not if $D_i$ is covered by $k-k'+1$ proper specializations missing an
instance of $\sigma_i$ that lies in no other target \src{untagged Prop 6.1}. For PA the gap between $k-c$ and $k-k'+1$ is
the whole story: $c=1$ for Sub-encoded steps, $c=4$ for raw sentences, and $c=7=k'-1$ with three negatives. The classical
sufficient condition ``not covered by $k$ failure sets'' is vacuous for $T_{\Ind}$ at $k\ge9$, because the nine root
specializations cover every induction instance. Plotkin's failure families are not $m$-noncovering for every $m$ once data
are read under universal closure (\cref{prop:many:heads}; the earlier contrary claim is withdrawn, \cref{app:ver:corr}).

\paragraph{Head classes.} In the rest of this subsection sentences are identified up to bijective renaming of parameters,
as the closure reading demands and as the implementation does (parameters are renamed canonically by first occurrence).
\Cref{rem:many:numerals} summarizes the results; \cref{prop:many:heads,prop:many:onevar,prop:many:unary} state and prove them.

\begin{remark}[head classes and numerals]\label{rem:many:numerals}
\status{proved; computed}\src{untagged U17}
Because data are read under universal closure, a one-variable instance schema $\varphi(z)$ is covered by its
$h=|F|+p+1$ head specializations ($F$ the function symbols, $p$ the number of parameters of $\varphi$; $5+p$ in
arithmetic), one per head class of the value of $z$, a parameter outside $\varphi$ being one class
(\cref{prop:many:heads}). So once $m\ge h$ (i.e.\ at least $h-1$ spare slots) the ``failure family'' covers everything.
For quantifier-free, parameter-free $\varphi$, under the hypotheses of \cref{prop:many:onevar}, the exact threshold is known
for $m\le h$ (\cref{prop:many:onevar}) and, over a unary signature, in closed form (\cref{prop:many:unary}). In particular
\emph{numeral-only data never make the cautious $k$-union learner accept $\forall x\varphi$ with one spare slot} when
$\varphi(a)$ lies in no other target's instance set: the sound split $\{\varphi(0)\}\cup\inst(\varphi(Su))$ misses
$\varphi(a)$ (\cref{thm:many:splits}(a)). Brute-force checks agree with all three propositions (2000/2000, 60/60, 254/254,
2509/2509 cases). Beyond this the spare-slot threshold is open (binary symbols with $m>h$, several metavariables,
templates with parameters, $T_{\Ind}$ with $m\ge9$, where the nine root specializations already cover $\inst(T_{\Ind})$).
\end{remark}

\begin{proposition}[finitely many head classes]\label{prop:many:heads}
\status{proved; computed}\src{untagged Prop 3.4}
Let $\sigma=\varphi(z)$ have a single metavariable, a 0-ary term metavariable $z$ occurring in $\varphi$; let $F$ be the
function symbols (constants included) and $c_1,\dots,c_p$ the parameters of $\varphi$. Every closed term $t$ has one head
class $\mathrm{head}(t)\in F\cup\{c_1,\dots,c_p\}\cup\{\nu\}$, where $\nu$ means ``a parameter other than the $c_j$''. Put
$\sigma_f:=\varphi(f(u_1,\dots,u_r))$ (fresh 0-ary $u$'s), $\sigma_{c_j}:=\varphi(c_j)$, $\sigma_\nu:=\varphi(b)$ for a
parameter $b$ not in $\varphi$. Then (a) $\inst(\sigma)=\bigcup_g\inst(\sigma_g)$ over the $h=|F|+p+1$ classes, and (b) no
$\sigma_g$ is redundant.
\end{proposition}
\begin{proof}
(a) For an instance $\varphi(t)$: if $t=f(\bar t)$ it is $\sigma_f[\bar u:=\bar t]$; if $t=c_j$ it is $\sigma_{c_j}$; if $t$ is
a parameter $b'\notin\{c_j\}$, the renaming swapping $b$ and $b'$ maps $\varphi(b)$ to $\varphi(b')$, since $\varphi$ contains
neither, so the two are the same datum. (b) Suppose a renaming $\pi$ maps $\varphi(t)$ to $\varphi(t')$. At the rigid
occurrences of $c_j$ in $\varphi$ it fixes $c_j$, so it maps other parameters to other parameters; at the position of $z$
it maps $t$ to $t'$. So $\mathrm{head}(t)=\mathrm{head}(t')$, and $\sigma_g$ covers only instances whose $z$-value has
class $g$.
\end{proof}
Computed\src{untagged u12(a)}: the 5 head-class templates of $z+0=z$ and the 6 of $c_1+z=z+c_1$ (one parameter) cover 2000/2000 random canonical
instances each, and each template is needed.

\begin{proposition}[threshold for one-variable instance schemas]\label{prop:many:onevar}
\status{proved; computed}\src{untagged Prop 3.5}
Assume (H1) $\sigma_i=\varphi(z)$ with $\varphi$ quantifier-free and parameter-free and $z$ a 0-ary term metavariable occurring
in it; (H2) every datum of $D_i$ has a $z$-value with at most one parameter, and $T$ is the set of these values with the
parameter renamed $a$; (H3) $N$ cross-separates $D$, the instance sets are pairwise disjoint, every $D_j$ is nonempty, and
$m=k-k'+1$. Let $\mathrm{cl}(T)$ be the set of head classes occurring in $T$ and $h=|F|+1$. (a) If $m<|\mathrm{cl}(T)|$,
then $\inst(\sigma_i)\subseteq\Acc_k(D,N)$. (b) If some class is missing from $T$, then $\inst(\sigma_i)\subseteq\Acc_k(D,N)$ iff
$m<|\mathrm{cl}(T)|$. (c) If all $h$ classes occur and $m=h$, then $\inst(\sigma_i)\subseteq\Acc_k(D,N)$ iff for each $f\in F$ of
arity $r\ge1$ the block $T_f$ of terms with head $f$ has $\lgg T_f=f(x_1,\dots,x_r)$ with distinct variables, i.e.\ Plotkin's
\evR{}+\evD{} holds for $T_f$ (each argument column has two different heads, the parameter counting as one head, and no two
columns are equal).
\end{proposition}

\begin{proof}[Proof of \cref{prop:many:onevar}]
\emph{Step 0 (least covering template of a block).} Let $\emptyset\neq B\subseteq D_i$ have $z$-values $\{t_b\}$. Every
$T\in\manyCov(B)$ satisfies $\inst(T)\supseteq\inst(\varphi(\lgg\{t_b\}))$, the lgg (Plotkin's) taken with $a$ as a constant.
Indeed: the data are quantifier-free, so $T$ has no binder in its rigid part, no bound variable is in scope, and every
metavariable is 0-ary: $T$ is a first-order pattern. It covers each datum through a representative
$\varphi(t_b[n_b/a])$. By Plotkin's theorem, $T\gen\lgg\{\varphi(t_b[n_b/a])\}=\varphi(G)$ with $G:=\lgg\{t_b[n_b/a]\}$ (lgg
commutes with the parameter-free context $\varphi$: rigid positions carry constant tuples, and the $z$-positions carry the
same tuple, hence the same variable). If all $n_b$ equal some $n$, $G=\lgg\{t_b\}[n/a]$, the same up to renaming.
Otherwise $G$ has no rigid parameter (a rigid $n$ at position $q$ would force $t_b|_q=a$ and $n_b=n$ for all $b$), and
$G\gen\lgg\{t_b\}$: the rigid symbols of $G$ are function symbols at positions where all $t_b$ agree, and two variable
positions of $G$ with the same tuple $(t_b[n_b/a]|_q)_b$ also have the same tuple $(t_b|_q)_b$ (undo the renamings), so
mapping each variable of $G$ to the corresponding subterm of $\lgg\{t_b\}$ is well defined and turns $G$ into
$\lgg\{t_b\}$. So $\inst(T)\supseteq\inst(\varphi(G))\supseteq\inst(\varphi(\lgg\{t_b\}))$. The template
$\varphi(\lgg\{t_b\})$ covers $B$, lies below $\sigma_i$, and so avoids every sound $N$.

\emph{Step 1 (reduction).} By \cref{thm:many:splits}(b), exactness for $i$ holds iff every $N$-avoiding family of at most
$m$ templates covering $D_i$ has union $\supseteq\inst(\sigma_i)$. Partitioning $D_i$ by first covering member and using
Step~0, this holds iff for every partition $\pi$ of $D_i$ (equivalently of $T$) into at most $m$ blocks,
$\bigcup_{B\in\pi}\inst(\varphi(\lgg T_B))=\inst(\sigma_i)$. Since $z$ occurs in the parameter-free $\varphi$, $\varphi(t)$
determines $t$ up to renaming, so this says that every closed term is, up to renaming, an instance of some $\lgg T_B$:
$\bigcup_{B\in\pi}\inst(\lgg T_B)=G_{\mathrm{cl}}$, the closed terms modulo renaming.

\emph{Step 2 (mixed blocks).} A block with two head classes has a variable as lgg, which covers $G_{\mathrm{cl}}$. So
failing partitions are head-pure, and a head-pure partition has at least $|\mathrm{cl}(T)|$ blocks. This proves (a).

\emph{Step 3 (a missing class).} If a class $g$ is missing from $T$, the head partition ($|\mathrm{cl}(T)|$ blocks) has lggs
that are constants, the parameter $a$, or terms with heads in $\mathrm{cl}(T)$, so no term of class $g$ is covered and it
fails. With (a), this proves (b).

\emph{Step 4 ($m=h$, all classes present).} The only head-pure partition with at most $h$ blocks is the head partition, and
it covers $G_{\mathrm{cl}}$ iff each block covers its class. A constant class (a 0-ary $f$, or $\nu$) has one term, which
its block covers. For $f$ of arity $r\ge1$, $\inst(\lgg T_f)\supseteq f(G_{\mathrm{cl}}^r)$ iff $\lgg T_f=f(x_1,\dots,x_r)$
with distinct variables: a non-variable argument misses terms with another head there, and a shared variable misses terms
with different subterms in the two arguments. Plotkin's construction gives a variable in argument $j$ iff column $j$ has
two different heads, and the same variable in two arguments iff the columns are equal.
\end{proof}

\begin{proposition}[closed form for unary signatures]\label{prop:many:unary}
\status{proved; computed}\src{untagged Prop 3.6}
Assume (H1)--(H3) with function symbols a finite set $C$ of constants and one unary $S$ (so (H2) is automatic and
$h=|C|+2$). Put $T_0:=T$, $T_{j+1}:=\{u:Su\in T_j\}$, and let $r$ be the least $j$ with $|\mathrm{cl}(T_j)|<h$. Then
$\inst(\sigma_i)\subseteq\Acc_k(D,N)$ iff $m<\kappa(T):=(h-1)r+|\mathrm{cl}(T_r)|$.
\end{proposition}

\begin{proof}[Proof of \cref{prop:many:unary}]
Let $\kappa(T)$ be the least number of blocks of a partition of $T$ whose lggs fail to cover $G_{\mathrm{cl}}$; by Steps
1--3, exactness holds iff $m<\kappa(T)$, and failing partitions are head-pure. If $r=0$ (a class is missing),
$\kappa(T)=|\mathrm{cl}(T_0)|$ by Step 3. If all $h$ classes are present, each of the $h-1$ constant classes ($C$ and
$\nu$) is a single term and forms a block of a head-pure partition, and the other blocks partition the $S$-class. Since
$\lgg$ of an $S$-block is $S$ applied to the lgg of the arguments and $\inst(Ss)=S(\inst s)$, the $S$-blocks cover
$S(G_{\mathrm{cl}})$ iff the corresponding blocks of $T_1$ cover $G_{\mathrm{cl}}$. So $\kappa(T)=(h-1)+\kappa(T_1)$, and
induction on $r$ gives $\kappa(T)=(h-1)r+|\mathrm{cl}(T_r)|$.
\end{proof}

Computed\src{untagged u12}: for $z+0=z$ and 60 random data sets, exactness at $m=h=5$, decided by brute force over all
partitions with coverage of all 608 canonical instances of depth $\le2$ (decisive, since every failure in Step~4 is
witnessed at depth $\le2$), agreed with (c) in 60/60 cases (18 exact), and with $m=|\mathrm{cl}(T)|-1$ the target was exact in
60/60; the closed form agreed with brute force on all 254 term sets of size $\le7$ over $\{S^jb:j\le3,b\in\{0,\nu\}\}$ and
all 2509 sets of size $\le6$ over $\{S^jb:j\le3,b\in\{0,c,\nu\}\}$.

%=====================================================================================================================
\subsection{Self-generated negatives}\label{app:many:pigeon}

\begin{corollary}[self-generated negatives]\label{cor:many:selfneg}
\status{proved}\src{untagged Cor 2.2}
Assume (Dec) and fix an effective enumeration of $\manyS$.
(a) If every $T\in\Min(\{a,b\})$ with $a\in A_i$, $b\in A_j$, $i<j$, is $d$-refuted, then one refuted instance $n_T$ per such
$T$ gives a set $N_\times$ satisfying (X) of \cref{thm:many:pigeonhole}, with $|N_\times|\le\mu\sum_{i<j}|A_i||A_j|$, where
$\mu=\max|\Min(\{a,b\})|$ over cross anchor pairs.
(b) The set $N(D,d)$ of first $d$-refuted instances of all $d$-refuted $T\in\Min(\{a,b\})$, $a,b\in D$, is computable from
$D$ without labels, is contained in $\Ref_d$, and contains a choice of $N_\times$. Hence $\Acc_{k'}(D,N(D,d))$ is sound at
all times, decidable, and equal to $\manyRstar$ as soon as $D$ contains anchors $A_i$ satisfying the hypothesis of (a).
\end{corollary}
\begin{proof}
(a) If $T'$ covers $a\in A_i$ and $b\in A_j$, then $T'\in\manyCov(\{a,b\})$, so $T'\gen T$ for some $T\in\Min(\{a,b\})$ and
$n_T\in\inst(T)\subseteq\inst(T')$. (b) $\Min(\{a,b\})$ is computable and ``$d$-refuted'' decidable; the first refuted
instance is found by enumeration. Within-target pairs may contribute further refuted sentences, which are false and
harmless. Apply \cref{lem:many:cautious} and \cref{thm:many:pigeonhole}.
\end{proof}

\begin{remark}[combining with the $k$-union learner]\label{rem:many:combine}
\status{proved}\src{untagged Rem 5.3}
Let $k\ge k'$ be known and accept $q$ iff $q\in\Acc_d\cap\Acc_k(D,N(D,d))$. This is target-sound at all times, by
\cref{lem:many:cautious}, with or without separation. Under $\RS_d(D)$ with $k=k'$, once every $D_i$ contains an anchor, it
equals $\manyRstar$: \DTRC{} is exact by \cref{thm:many:dtrc}(iii), and $\RS_d(D)$ makes every minimal template of every
cross pair $d$-refuted, so $N(D,d)$ satisfies (X) and \cref{thm:many:pigeonhole} applies. The anchor hypothesis is
needed for this argument: \DTRC{} can be exact without anchors, because refuted templates of whole clusters drop out,
while $N(D,d)$ contains only refutations of pair templates (\cref{app:ver:corr}).
\end{remark}

For PA (\cref{tab:many:pacross}) the four sentences $0=S0$, $0=0\to0=S0$, $p+0=0$, $p\cdot0=S0$ refute every cross template;
since the $\Q$/induction templates $F_0$ and $F_0\to F_1$ contain all of $\inst(T_{\Ind})$, three suffice for the
$k$-union learner (\cref{cor:many:paslots}).

%=====================================================================================================================
\subsection{\texorpdfstring{\DTRC}{DTRC}: proofs}\label{app:many:dtrc}

\begin{proof}[Proof of \cref{thm:many:dtrc}]
(i) By \cref{prop:many:W}(c) some $T_0\in\Min(X)$ has $\inst(T_0)\subseteq\inst(\sigma_i)\subseteq\manyRstar\subseteq\manyTr$,
so $\inst(T_0)\cap\Ref_d=\emptyset$ by (OS).
(ii) Call a cluster pure if it lies inside one $D_i$; initially all are. If $C\subseteq D_i$ and $C'\subseteq D_j$ are pure
with $i\neq j$, every $T\in\manyCov(C\cup C')$ covers a datum of $D_i$ and one of $D_j$, so it is $d$-refuted by $\RS_d(D)$:
the union is not coherent and is never formed. So all clusters stay pure. The process stops after finitely many merges.
At termination, two final clusters inside one $D_i$ would have a coherent union by (i), contradicting termination. So
each nonempty $D_i$ is exactly one final cluster, whatever the policy.
(iii) For the final cluster $D_i$, the $T_0$ of (i) lies in $\Min_d^+(D_i)$, so $\Acc_d(D_i)\subseteq\inst(T_0)\subseteq
\inst(\sigma_i)$ and $\Min^+_d(D_i)\neq\emptyset$. If $A_i\subseteq D_i$ is an anchor, every $T\in\Min(D_i)$ covers $A_i$, so
$\inst(T)\supseteq\inst(\sigma_i)$ and $\Acc_d(D_i)\supseteq\inst(\sigma_i)$.
(iv) Under global separation, (ii) and (iii) hold for every $D\subseteq\manyRstar$. So \DTRC{} fails to be exact only if some
$D_i$ lacks an anchor, which is the failure event of the tagged learner given the labels; the union bound over $i$ gives
the sum, and averaging the tagged bound over the binomial number of target-$i$ data gives the rate.
(v) A pair answer depends only on the two clusters. There are $n(n-1)/2$ initial pairs; there are at most $n-1$ merges,
and after each only pairs with the new cluster (at most $n$) need retesting. For the online policy under $\RS_d(D)$, by
induction on the data there is at most one cluster per target at all times: a datum of target $i$ joins the first coherent
cluster, cross unions are incoherent, and the cluster of target $i$ is coherent with it by (i). So each datum is tested
against at most $k'$ clusters, and the final pass tests the $k'(k'-1)/2$ pairs of pure clusters.
\end{proof}

\begin{proof}[Proof of \cref{thm:many:residue}]
(a) A pure final cluster $C\subseteq D_i$ contributes $\Acc_d(C)\subseteq\inst(\sigma_i)$, as in the proof of
\cref{thm:many:dtrc}(iii). A mixed final cluster $C$ is $d$-coherent or an incoherent singleton contributing $\emptyset$
(\cref{prop:many:linkage}), and it contributes $\Acc_d(C)\subseteq\inst(T)$ for each $T\in\Min_d^+(C)$, which is part of
$\Res_d(D)$ (take $X=C$). (b) $\Acc_d(C)\subseteq\inst(T)\subseteq\manyTr$. (c) By definition; $\Res_d(D)$ decreases in $d$
because coherence and non-refutation are antitone in $d$ (\cref{lem:many:mono}). (d) In the example of \cref{sec:many:residue}
all data are quantifier-free, so every minimal template is a Plotkin lgg (\cref{prop:many:testcost}(b)). $\{g,0+0=0\}$ has
the lgg $0+z=z$ (the columns $(S0,0)$ at the second argument and at the right-hand side are equal), which is true and never
refuted; $\{g,0+0=0,S0+0=S0\}$ has $z_0+z_1=z_2$ and $\{S0+0=S0,g\}$ has $z_0+z_1=S0$, both refuted by $0+0=S0$. So in the
order $g,0+0=0,S0+0=S0$ the final clusters are $\{g,0+0=0\}$ and $\{S0+0=S0\}$; $p+0=p$ is in neither acceptance set,
while $0+p=p$ is accepted. Computed\src{untagged u9, Example B}, with the referee's variant with overlapping instance sets
($\sigma_2=0+z=z$) behaving the same way in 2 of 3 orders.
\end{proof}

\begin{proof}[Proof of \cref{lem:many:twopoint}]
\emph{Two-point identity.} Let $G(z),K(z)$ be expressions in which $z$ occurs only as a leaf, with $G(t_j)=K(t_j)$ for
$j=1,2$. Then $G=K$. Induction on $G$ and $K$: if $G=z$ and $K\neq z$, then $K$ has a rigid head $h$ (a symbol or an
index) that survives substitution, and $K(t_j)=t_j$ forces $h=\mathrm{head}(t_1)$ and $h=\mathrm{head}(t_2)$, a
contradiction; symmetrically for $K=z$. If neither is $z$, their heads agree (compare at $t_1$), and we recurse into the
arguments; binders are harmless, because $z$ is not a bound variable and the $t_j$ are closed.
\emph{Skeleton.} The two data agree everywhere except at and below the $z$-positions of $\varphi$, where their heads
differ. By \cref{lem:single:basic}(c) the rigid skeleton of a covering template $T$ lies in the common
prefix of the data, so it avoids the $z$-positions and everything below them; every occurrence of $T$ is at a position of
$\varphi$.
\emph{Bodies.} For each metavariable $M$ of $T$ choose a pattern occurrence $M(\bar y)$ at position $p$. For any closed
$t$ put $\theta_t(M):=\lambda\bar y.(\varphi(t)|_p)$. This is a legal body: $t$ is closed, so the bound variables free in
$\varphi(t)|_p$ are those free in $\varphi(t_1)|_p$, which are among $\bar y$ because $T$ covers $\varphi(t_1)$.
\emph{Every occurrence fits.} For an occurrence $M(\bar u)$ at position $p'$ put $G(z):=(\lambda\bar y.\varphi(z)|_p)(\bar u)$
(plugging) and $K(z):=\varphi(z)|_{p'}$. $G$ contains $z$ only as a leaf, since $\bar u$ is metavariable-free. $T$ covers
both data with the unique matchers $\theta_{t_1},\theta_{t_2}$, and substituting a closed $t_j$ for $z$ commutes with
abstraction and plugging, so $G(t_j)=K(t_j)$; hence $G=K$ and $G(t)=K(t)$. With the skeleton, $T\theta_t=\varphi(t)$, so
$\inst(T)\supseteq\inst(\varphi(z))$.
\end{proof}

\begin{proof}[Proof of \cref{prop:many:forall}]
(a) $\varphi(z)$ covers $D$, and by \cref{lem:many:twopoint} every covering template contains $\inst(\varphi(z))$; so some
covering template is unrefuted iff $\varphi(z)$ is (\cref{lem:many:mono}). (b) In closure-normal form a parameter is read
universally; it must be fresh, so that $\varphi(a)$ is not an instance at a special parameter. If all instances are true,
so is the closure of $\varphi(a)$, i.e.\ $\forall x\varphi$; conversely every instance follows from $\forall x\varphi$.
(c) If $\forall x\varphi$ is true, every instance is true, so every subset of $D$ is coherent (\cref{thm:many:dtrc}(i) with
the single target $\varphi(z)$), and no two clusters can remain at termination. In the single final cluster every
unrefuted minimal template contains $\inst(\varphi(z))$ (\cref{lem:many:twopoint}), and one of them is contained in it
(\cref{prop:many:W}(c)), so $\Acc_d=\inst(\varphi(z))$. The oracle is a fixed function of the query, so the practices
``$n$ ground instances'' and ``$\varphi(z)$'' produce the same run. (d) By (a); a subset whose terms all have the same head,
e.g.\ $S$, has other minimal templates such as $\varphi(Sz)$, which need not be refuted. (e) The first claim is the
definition of completeness for false instances plus (b), and \cref{prop:many:qf} applied to $T=\varphi(z)$ gives it for
R-$\Delta_0$ and quantifier-free $\varphi$. R-$\Delta_0$ verifies universal statements only by equality reasoning from true
closed equations, so a false sentence whose refutation needs a true universal premise that this reasoning cannot verify is
never refuted; the noisy-PA computation (\cref{ex:many:noise}) contains such a sentence, the wrong-base induction variant
$(\neg S0=0\wedge\forall x(\neg x=0\to\neg Sx=0))\to\forall x\,\neg x=0$, whose refutation needs the universal premise
$\forall x(\neg x=0\to\neg Sx=0)$, which follows from Q1 \src{untagged referee U7; u8}. At a finite budget, a merge through
a false universal whose least counterexample lies beyond the budget is part of $\Res_d$ (\cref{thm:many:residue}). In $\R$
as an ordered field every closed term denotes an integer (\cref{ex:univ:fail}(a)), and no integer squares to $2$, so every
closed instance of $\neg(z\cdot z=1+1)$ is true, while $\neg(p\cdot p=1+1)$, read universally, is false; a decision
procedure for real closed fields would refute it. $\Q\vdash0+\num n=\num n$ for each $n$ but $\Q\nvdash\forall x(0+x=x)$ is
\cref{ex:univ:q}.
\end{proof}

\noindent Computed\src{untagged u6}: numeral instances of Q4--Q6, as 12 ground targets, form three clusters with templates
$z_0+0=z_0$, $z_0+Sz_1=S(z_0+z_1)$, $z_0\cdot0=0$, and \DTRC{} accepts $p+0=p$; instances of the false $n\cdot n=n$
($n=0,1$) stay apart (refuted at $2$); $\neg(n\cdot n=2)$, $n\le2$, merge and the true $\neg(p\cdot p=2)$ is accepted.

\begin{proof}[Proof of \cref{prop:many:qf}]
Let $s=T\theta$ be false. Its universal closure is false in $\N$, so some assignment $\bar n$ of numerals to the
parameters of $s$ makes $s[\bar n]$ false. Let $\theta'$ agree with $\theta$ on term metavariables, and for a formula
metavariable $F$ let $\theta'(F):=(0=0)$ if $\theta(F)[\bar n]$ is true in $\N$ and $(0=S0)$ otherwise. Then $s':=T\theta'$
is quantifier-free, since $T$'s rigid part, the term values and the new formula values are. Moreover $s'[\bar n]$ has the
truth value of $s[\bar n]$: the rigid part has no binders, so it is a Boolean combination of equations between terms built
from rigid symbols and unchanged term values, and of occurrences of formula metavariables, whose replacements have the
same truth values at $\bar n$. So $s'$ is false. At depth at least $\max(\bar n,|s'|)$, R-$\Delta_0$ instantiates the
parameters of $s'$ (which are among those of $s$) by $\bar n$ and evaluates the resulting closed quantifier-free sentence,
so $s'\in\Ref_\infty$. If $T$ is in $\Res_\infty$, no instance is refuted, so $T$ has no false instance.
\end{proof}

\begin{proof}[Proof of \cref{thm:many:depth}]
\emph{Construction.} By the MRDP theorem \citep{davis1961decision,matiyasevich1993hilbert} there are polynomials
$P(e,\bar x)$, $Q(e,\bar x)$ with natural coefficients, $\bar x=x_1\dots x_k$, such that $e\in K\iff\exists\bar x\in\N^k\,
P(e,\bar x)=Q(e,\bar x)$, $K$ the halting set. Let $P_e,Q_e$ be the terms with the numeral $\num e$ substituted, and put
\[
H_e(z):=\exists x_1\dots\exists x_k\,(z=SS(x_1+\dots+x_k)\wedge P_e(\bar x)=Q_e(\bar x)),
\]
with a 0-ary term metavariable $z$ under $k$ binders ($z$ is closed, so no capture). Let $T_e:=\neg H_e(z)$, $A_e:=\neg H_e(0)$,
$B_e:=\neg H_e(S0)$, $q_e:=\neg H_e(SS0)$. Practice $\mathcal P_e$ has the single target $T_e$; practice $\mathcal P'_e$ has the ground targets
$A_e,B_e$; both have the data law uniform on $\{A_e,B_e\}$. \emph{Facts.} (1) $A_e,B_e$ are true, since $0\neq SSu$ and
$S0\neq SSu$; so $\mathcal P'_e$ is valid. (2) $\mathcal P_e$ is valid iff $e\notin K$: its instances, including the parameter instance read
universally, are all true iff no solution exists. (3) $\{A_e,B_e\}$ is an anchor for $T_e$ (\cref{lem:many:twopoint}: heads $0$
and $S$). (4) If $e\in K$ with solution $\bar x^*$, let $\bar c$ be the numeral of $2+\sum x_i^*$. The instance $\neg H_e(\bar
c)$ is false, and R-$\Delta_0$ refutes it at depth $d_e:=\max(|\neg H_e(\bar c)|,\max_ix^*_i)$: refuting $\neg\psi$ means
verifying $\psi=\exists\bar x(\dots)$, done by the witnesses $\bar x^*$ and closed evaluation of two atoms; no universal has
to be verified. By (3) every template covering $A_e$ and $B_e$ contains $\neg H_e(\bar c)$, so $\mathcal P'_e$ is separated at depth
$d_e$ and $\Res_\infty(\mathcal P'_e)=\emptyset$. (5) If $e\notin K$, $T_e$ is truth-sound and never refuted, and
$\Res_\infty(\mathcal P'_e)\supseteq\inst(T_e)\ni q_e$; $\mathcal P'_e$ is separated at no depth, so (b) asks nothing of it. (6) If $e\notin K$,
$\mathcal P_e$ has a single target, is separated vacuously, and its data contain the anchor $\{A_e,B_e\}$ with positive
probability.
\emph{Reduction.} Let $S:=\{e:\exists t\ \Pr[q_e\in\Acc_{t'}\text{ for some }t'\le t]>\frac12\}$. If $e\notin K$, (b) for
$\mathcal P_e$ gives eventual acceptance of $q_e\in\manyRstar_{\mathcal P_e}$ with probability $\ge1-\delta>\frac12$; by continuity along the
increasing events ``accepted by time $t$'', some single $t$ has probability $>\frac12$, so $e\in S$. If $e\in K$, (a) for
$\mathcal P'_e$ (with $\manyRstar=\{A_e,B_e\}$, $\Res_\infty=\emptyset$, $q_e\notin\manyRstar$) gives $\Pr[q_e\text{ ever accepted}]\le
\delta<\frac12$, so $e\notin S$. Hence $S=\omega\setminus K$. The learner, the data law and the oracle are computable
uniformly in $e$ ($\Ref_d$ consists of the sentences of size $\le d$ refuted with numeral bounds $d$, decidable uniformly in
$d$), and $\mathcal P_e$, $\mathcal P'_e$ have the same data law and oracle; so the probability that $q_e$ is accepted by time $t$ is
computable from $(e,t)$ (for randomized learners, approximable from below), and $S$ is $\Sigma_1$. But $\omega\setminus K$
is not $\Sigma_1$.
\end{proof}
R-$\Delta_0$ cannot verify the bounded universal quantifiers of Kleene's $T$ predicate, which is why the construction
uses MRDP (\cref{app:ver:corr}). Alternatively, keep the Kleene-$T$ form and take as oracle the full evaluation of
closed $\Delta_0$ sentences with bounded quantifiers as primitives, which is also sound and decidable per depth.

%=====================================================================================================================
\subsection{Whole-cluster tests and noise}\label{app:many:noise}

\begin{proposition}[whole-cluster versus pairwise tests]\label{prop:many:linkage}
\status{proved; computed}\src{untagged Prop 5.8}
Coherence is downward closed but not transitive. Every cluster built by \DTRC{} is coherent at all times (except a
refutable singleton, which contributes nothing), whereas single linkage on the pairwise-coherence graph can build clusters
with no unrefuted minimal template. Under $\RS_d(D)$, the coherent subsets of $D$ are exactly the subsets of the $D_i$, and
single linkage, complete linkage and \DTRC{} agree. Without separation, \DTRC{} returns some maximal partition into
coherent blocks, which may depend on the policy.
\end{proposition}
Example: $a=(0+0=0)$, $b=(S0+0=S0)$, $c=(0+S0=S0)$. The pairs $\{a,b\}$, $\{a,c\}$ have the unrefuted templates
$z_0+0=z_0$, $0+z_0=z_0$, while $\{b,c\}$ and $\{a,b,c\}$ have only $z_0+z_1=S0$, resp.\ $z_0+z_1=z_2$, refuted by
$0+0=S0$. Single linkage lumps all three; \DTRC{} returns $\{\{a,b\},\{c\}\}$ or $\{\{a,c\},\{b\}\}$.

\paragraph{Noise.} Let $D=D_{\mathrm{clean}}\cup E$ with mistakes $E\cap\manyRstar=\emptyset$.

\begin{proposition}[noise]\label{prop:many:noise}
\status{proved; computed}\src{untagged Prop 5.9}
(a) A refutable mistake stays an incoherent singleton. (b) Under $\RS_d(D_{\mathrm{clean}})$ no coherent set contains clean
data of two targets, so mistakes never bridge targets. (c) An unrefutable mistake can be absorbed into a target's cluster
and make it unsound, and (d) it can \emph{fragment} the target: $C\cup\{m\}\cup C'$ can be incoherent although $C\cup C'$
is coherent. (e) For $|C|>e$ the \emph{trimmed verifier} $\Acc_d^e(C):=\bigcap\{\inst(T):T\text{ not }d\text{-refuted},\
|C\setminus\inst(T)|\le e\}$ is contained in $\inst(\sigma_i)$ when $C$ has at most $e$ mistakes and its clean part lies in
$D_i$, and equals $\inst(\sigma_i)$ when moreover every subset of $C$ of size $\ge|C|-e$ contains an anchor.
\end{proposition}
\emph{Robust} \DTRC{} asserts only clusters of multiplicity $\ge s>e$, verified with $\Acc^e_d$. If $\RS_d(D_{\mathrm{clean}})$
holds, every mistake is refutable or \emph{noise-separated} (every template covering it and a clean datum is refuted), and
every coherent set of mistakes has multiplicity $<s$, then the asserted clusters are exactly the $D_i$ of multiplicity
$\ge s$, verified soundly, and exactly once every subset of $D_i$ of size $\ge|D_i|-e$ contains an anchor. Fragmentation
is not prevented, and a frequent coherent family of mistakes is a systematic error that no positive-data method can tell
from a rule. A computed example on PA (\cref{ex:many:noise}) shows absorption of the false instance
$J^*_0=(0+0=0\wedge\forall x(0+x=0\to0+Sx=S0))\to\forall x(0+x=0)$ into a non-diverse induction cluster, its rejection by
the trimmed verifier, and order-dependent fragmentation.

\begin{proof}[Proof of \cref{prop:many:linkage}]
Non-transitivity and policy dependence: the example after the proposition (all data quantifier-free, so the minimal
templates are Plotkin lggs). \DTRC{} merges only coherent unions, so its clusters stay coherent, apart from a singleton
$\{m\}$ with $m\in\Ref_d$, which no merge involves since every covering template of a set containing $m$ has $m$ as an
instance. Under $\RS_d(D)$ the subsets of the $D_i$ are coherent (\cref{thm:many:dtrc}(i)) and every other set is not, so the
pairwise-coherence graph is a disjoint union of cliques $D_i$, and all three procedures return $\{D_i\}$.
\end{proof}

\begin{proof}[Proof of \cref{prop:many:noise}]
(a) $m\in\inst(T)$ for every $T\in\manyCov(X)$ with $m\in X$. (b) A template covering clean $a\in D_i$ and $b\in D_j$ covers
$\{a,b\}$ and is $d$-refuted. (c) By definition; \cref{ex:many:noise} gives an instance. (d) \Cref{ex:many:noise}.
(e) First, the definition of $\Acc^e_d(C)$ agrees with the finite form
\[\Acc^e_d(C)=\bigcap_{E'\subseteq C,\,|E'|\le e}\ \bigcap_{T\in\Min^+_d(C\setminus E')}\inst(T):\]
every unrefuted $T$ with $|C\setminus\inst(T)|\le e$ covers $C\setminus E'$ for $E':=C\setminus\inst(T)$ and lies above some
$T_0\in\Min(C\setminus E')$, which is unrefuted; conversely each such $T_0$ is an unrefuted template missing at most $e$
data ($C\setminus E'\neq\emptyset$ since $|C|>e$). Soundness: take $E'$ to be the mistakes; $C\setminus E'\subseteq D_i$, so
\cref{prop:many:W}(c) gives an unrefuted $T_0\in\Min(C\setminus E')$ with $\inst(T_0)\subseteq\inst(\sigma_i)$. Exactness: for
each $E'$, $C\setminus E'$ contains an anchor, so every member of $\Min(C\setminus E')$ contains $\inst(\sigma_i)$.
\emph{Robust} \textup{\DTRC}. Refutable and noise-separated mistakes never join a cluster with clean data, so by (b) and
\cref{thm:many:dtrc}(ii) the clean data form the clusters $D_i$, and clusters of mistakes have multiplicity $<s$ and are not
asserted. Each asserted $D_i$ has no mistakes, so (e) applies with zero mistakes. If noise-separation is dropped, each
asserted cluster still has clean data of one target (by (b)) and, by assumption, at most $e$ mistakes, so (e) gives
soundness and the stated exactness per cluster.
\end{proof}

\begin{example}[noisy PA]\label{ex:many:noise}
\status{computed}\src{untagged u8; reproduced by the referee}
To 60 clean PA data add three refutable mistakes (e.g.\ $p+0=Sp$), two unrefutable false ones ($J^*_0$, and the wrong-base
variant $(\neg S0=0\wedge\forall x(\neg x=0\to\neg Sx=0))\to\forall x\,\neg x=0$, whose refutation needs Q1), and two true
non-targets ($0+p=p$, $\neg Sp=p$). The refutable mistakes stay incoherent singletons; the other mistakes stay singletons of
multiplicity 1, since their merges with the diverse induction cluster are refuted, so with $s=2$ none is asserted.
\emph{Absorption}: $J^*_0$ together with the non-diverse data $\Ind(S^n0+x=S^nx)$, $n<3$, is coherent, with the single
unrefuted minimal template $(f(0)+0=f(0)\wedge\forall x(f(0)+x=f(x)\to f(0)+Sx=Sf(x)))\to\forall x(f(0)+x=f(x))$, which is
$K_0$-shaped and not below $T_{\Ind}$ (unique among 999 covering templates of size $\le20$, by the referee's enumeration).
The plain verifier accepts the false $J^*_0$ and $J^*_1=(S0+0=S0\wedge\forall x(S0+x=S0\to S0+Sx=SS0))\to\forall x(S0+x=S0)$;
the trimmed verifier with $e=1$ rejects both and accepts the genuine $\Ind(S^30+x=S^3x)$. \emph{Fragmentation}: in the order
$\Ind(0+x=x)$, $J^*_0$, $\Ind(S0+x=Sx)$, then three diverse inductions, \DTRC{} returns the unsound cluster
$\{\Ind(0+x=x),J^*_0,\Ind(S0+x=Sx)\}$ and a cluster of the diverse inductions; with the diverse data first it puts all five
inductions together and leaves $J^*_0$ alone. Merging most specific or largest clusters first is a heuristic remedy, not a
theorem.
\end{example}

%=====================================================================================================================
\subsection{Complexity of untagged verification and of the refutation test}\label{app:many:cost}

\begin{proposition}[cost of a test]\label{prop:many:testcost}
\status{proved; computed}\src{untagged Prop 5.10}
(a) Given $\Min(X)$, a $d$-coherence test needs at most $|\Min(X)|$ refutation checks. (b) If $X$ is quantifier-free, every
covering template is a first-order pattern and $\Min(X)=\{\lgg(X)\}$ (Plotkin's lgg with 0-ary term and formula variables),
computable in polynomial time. (c) If $X\subseteq\inst(\sigma)$ contains an anchor of $\sigma$, then, under \Rich,
$\Min(X)=\{\sigma\}$, computable from $X$ in polynomial time. (d) For one negative $n$: some $T\in\Min(X)$ avoids $n$ iff
$n\notin\Acc(X)$, decidable in polynomial time. (e) In every coherence test of the PA runs ($63+102+98$ tests) and ZF runs
($164+131+147$ tests) of \cref{sec:many:examples}, $|\Min(X)|=1$; the worst case met was in low-diversity induction pairs
($|\Min|$ up to 27).
\end{proposition}
\begin{proof}
(a) By definition. (b) The data have no binders, so neither has the rigid part of a covering template; no bound variable
is in scope, so pattern occurrences, and hence all metavariables, are 0-ary. Plotkin's lgg is least among first-order
patterns \citep{plotkin1970note,reynolds1970transformational}. (c) $X$ is itself an anchor (its covering templates cover
the anchor), so \cref{cor:single:anchor-lgg} makes $\sigma$ the least covering template, computed by the polynomial lgg
construction of \cref{prop:single:lgg}. (d) If $n\notin\Acc(X)=\bigcap\manyCov(X)$, some covering $T$ misses
$n$, and so does a minimal $T_0$ below it; the converse is trivial. $\Acc(X)$ is decided in polynomial time by the feature
verifier (\cref{thm:single:feat,cor:single:poly}). (e) Computed\src{untagged u3, u5, u0}.
\end{proof}

\paragraph{Features.} Common prefix $C(B)$, slots, $A(B)$, scopes $Y_\sigma(B)$, the features $\mathrm{Sym}(p)$,
$\mathrm{Scope}(\sigma)$, $\mathrm{Eq}(\sigma,r,\bar u)$ and $\Feat(B)$ are as in \cref{def:single:features}, and
$\Acc(B)=\Feat(B)$, decidable in polynomial time (\cref{thm:single:feat,cor:single:poly}). For positions $\sigma\perp r$ of the
same sort and a sentence $d$, the \emph{matcher} $e_d$, if it exists, is the unique partial map on $\FV(d|_\sigma)$ with
$d|_r=(d|_\sigma)[e_d]$ (forcing, \cref{lem:single:basic}(d)). With \eqref{eq:many:partition}, ``$q\in\Acc_k(D)$'' is in
coNP for every $k$: a partition into at most $k$ blocks $B$ with $q\notin\Acc(B)$ is a certificate of non-membership.

\begin{lemma}[failure types]\label{lem:many:types}
\status{proved; computed}\src{single Lemma H.1}
Fix $q$. Say that $d$ \emph{agrees with $q$ above $p$} if every strict ancestor of $p$ is a position of $d$ with the same
root symbol as in $q$ (then $p$ is a position of $d$). For nonempty finite $B$, $q\notin\Acc(B)$ iff $B$ fits one of the
following \emph{types}, indexed by positions of $q$:
\begin{itemize}
\item $\mathrm{Sym}(p,h)$, $h\neq\mathrm{root}(q|_p)$: every $d\in B$ agrees with $q$ above $p$ and has root $h$ at $p$;
\item $\mathrm{Scope}(\sigma,y)$, $y\in\FV(q|_\sigma)$: every $d\in B$ agrees with $q$ above $\sigma$, and
  $y\notin\FV(d|_\sigma)$;
\item $\mathrm{Eq}(\sigma,r)$, $\sigma\perp r$ of the same sort: every $d\in B$ agrees with $q$ above $\sigma$ and above $r$
  and has a matcher $e_d$; the $e_d$ are pairwise compatible; and either $e_q$ does not exist or it is incompatible with
  some $e_d$, $d\in B$.
\end{itemize}
There are at most $|q|(|\Sigma_D|+2b_{\max}(q))+|q|^2$ types ($\Sigma_D$ the non-index symbols of the data, $b_{\max}(q)$ the
binder depth of $q$). For each type, ``$B$ fits'' is a conjunction of: $B\subseteq$ a fixed \emph{allowed} set; $B$ a clique of
a fixed compatibility graph (Eq-types); $B$ meets a fixed \emph{witness} set (only for Eq-types where $e_q$ exists).
\end{lemma}
\begin{proof}
($\Rightarrow$) $q$ lacks some $B$-feature; take the first applicable case. If $q$ lacks $\mathrm{Sym}(p)$ for some $p\in
C(B)$, let $p'$ be the topmost position on the path to $p$ where $q$'s root differs from $B$'s common root $h'$; its strict
ancestors agree with $q$, so $p'$ is a position of $q$, and $B$ fits $\mathrm{Sym}(p',h')$. If $q$ has every Sym feature but
lacks $\mathrm{Scope}(\sigma)$, the strict ancestors of $\sigma$ lie in $C(B)$ with $q$'s symbols; any $y\in\FV(q|_\sigma)
\setminus Y_\sigma(B)$ gives the type $\mathrm{Scope}(\sigma,y)$. If $q$ has every Sym and Scope feature but lacks
$\mathrm{Eq}(\sigma,r,\bar u)$, each $e_d$ exists and, by forcing, is the restriction of $\bar u$ to $\FV(d|_\sigma)$, so the
$e_d$ are pairwise compatible; also $\FV(q|_\sigma)\subseteq Y_\sigma(B)=\mathrm{dom}\,\bar u$ and $q|_r\neq(q|_\sigma)[\bar u]$. So
either $e_q$ does not exist, or $e_q(y)\neq\bar u(y)=e_d(y)$ for some $y$ and some $d$ with $y\in\FV(d|_\sigma)$.

($\Leftarrow$) By induction on the size of $q|_\sigma$ (resp.\ $q|_p$). The named positions lie in $A(B)$, since their
strict ancestors lie in $C(B)$ with $q$'s symbols. \emph{$\mathrm{Sym}(p,h)$}: $p\in C(B)$, and $q$ lacks $\mathrm{Sym}(p)$.
\emph{$\mathrm{Scope}(\sigma,y)$}: if $\sigma$ is a slot, $q$ lacks $\mathrm{Scope}(\sigma)$. If $\sigma\in C(B)$ with common
root $g$ and $q$'s root differs, $\mathrm{Sym}(\sigma)$ fails. Otherwise $g$ is not an index ($g=y$ would put $y$ into every
$\FV(d|_\sigma)$; another index would give $\FV(q|_\sigma)=\{g\}\not\ni y$), so $y$ occurs, shifted if $g$ binds, in some child
$q|_{\sigma\cdot i}$, and $B$ fits $\mathrm{Scope}(\sigma\cdot i,y')$; apply induction. \emph{$\mathrm{Eq}(\sigma,r)$}: if
$\sigma$ is a slot, $\bar u:=\bigcup_de_d$ is well defined on $Y_\sigma(B)$ by compatibility, $\mathrm{Eq}(\sigma,r,\bar u)$ is a
$B$-feature, and $q$ has it only if $e_q$ exists and $e_q\subseteq\bar u$, which the type excludes. If $\sigma\in C(B)$ with
common root $g$ and $q$'s root differs, $\mathrm{Sym}(\sigma)$ fails. If $g$ is an index (free at $\sigma$), then
$d|_r=e_d(g)$ for each $d$, and compatibility makes all $d|_r$ equal to one term $u$, so $r$ and the positions of $u$ below it
lie in $C(B)$; $e_q$ exists with $e_q(g)=q|_r$, and the type forces $q|_r\neq u$, so a Sym feature fails inside $u$. If $g$
is a non-index symbol, every $d|_r=(d|_\sigma)[e_d]$ has root $g$, so $r\in C(B)$, and if $q|_r$ has another root,
$\mathrm{Sym}(r)$ fails. Otherwise pass to the children: the restricted matchers $e^{(i)}_d$ exist and are pairwise
compatible (a local index bound by $g$ is mapped to itself). The failure of $e_q$ shows at some child $i$ and some index
$z$ free in $q|_{\sigma\cdot i}$: $q$'s child matcher does not exist, maps $z$ to a wrong value, or two children of $q$
disagree on $z$. If some $d\in B$ has $z$ free in $d|_{\sigma\cdot i}$, compatibility fixes the correct value and $B$ fits
$\mathrm{Eq}(\sigma\cdot i,r\cdot i)$ (in the disagreement case, at the child with the wrong value); otherwise $B$ fits
$\mathrm{Scope}(\sigma\cdot i,z)$. Apply induction.

The count: Sym-types are indexed by a position of $q$ and a symbol of the data or one of fewer than $b(p)$ indices,
Scope-types by a position and an index, Eq-types by a pair of positions. The conditions are read off the definitions.
\end{proof}
Computed\src{single e11}: on 600 random cases (mixtures of 1--3 random $\DT$ templates with low-diversity bodies, 3--6 data),
``$q\notin\Acc(B)$'' agreed with ``$B$ fits a type'' on all 118{,}818 pairs $(B,q)$, with 15{,}220 Eq-types having
non-trivial compatibility graphs.

\begin{proof}[Proof of \cref{thm:many:complexity}(a)]
$k=1$ is \cref{thm:single:feat}. For $k=2$: by \eqref{eq:many:partition}, $q\notin\Acc_2(D)$ iff $q\notin\Acc(D)$ or there are
types $\tau_1,\tau_2$ and a partition $D=B_1\sqcup B_2$ with $B_l$ fitting $\tau_l$. Fix $(\tau_1,\tau_2)$ and, for Eq-types that
need one, witnesses $w_l\in B_l$ (polynomially many tuples). With variables $x_d$ (``$d\in B_1$''), ``$B_l$ fits $\tau_l$'' is
a conjunction of clauses with at most two literals: $\neg x_d$ for $d$ not allowed by $\tau_1$; $x_d$ for $d$ not allowed by
$\tau_2$; $\neg x_d\vee\neg x_{d'}$ for $d,d'$ incompatible for $\tau_1$; $x_d\vee x_{d'}$ for $d,d'$ incompatible for $\tau_2$;
and $x_{w_1}$, $\neg x_{w_2}$. 2-SAT is solvable in linear time \citep{aspvall1979linear}. The algorithm first checks $q\in
\Acc(D)$; after that, a solution with an empty block would make $D$ fit a type, i.e.\ $q\notin\Acc(D)$, so every solution is a
genuine 2-partition.
\end{proof}

\begin{proof}[Proof of \cref{thm:many:complexity}(b)]
Membership in coNP was noted above. Hardness: reduction from graph $k$-colourability, NP-complete for each fixed $k\ge3$
\citep{karp1972reducibility,garey1979computers}. Let $G=(V,E)$ with $V=\{1,\dots,n\}$, $E=\{e_1,\dots,e_m\}$, each edge
oriented (tail, head). Under $m$ binders $\forall y_1\dots\forall y_m$ put $\mathrm{tup}_u:=t_{u,1}+(t_{u,2}+(\dots+(t_{u,m}+0)))$
with $t_{u,i}=y_i$ if $u\in e_i$ and $0$ otherwise; let $c_u$ map $y_i$ to $0$ if $u$ is the tail of $e_i$ and to $S0$ if it is
its head; and put
\[
d_u:=\forall\bar y\,\bigl(S^u(\mathrm{tup}_u)=0\wedge S^u(\mathrm{tup}_u[c_u])=0\bigr),\qquad
q:=\forall\bar y\,\bigl(S^n(0)=0\wedge S^{n+1}(0)=0\bigr),
\]
of total size $O(n(n+m))$. Let $\sigma_w,r_w$ be the positions at $S$-depth $w$ in the left-hand sides of the two conjuncts.
\emph{Claim: for nonempty $B\subseteq D$, $q\notin\Acc(B)$ iff the vertices of $B$ form an independent set.} Then by
\eqref{eq:many:partition}, $q\notin\Acc_k(D)$ iff $V$ splits into at most $k$ independent sets, iff $G$ is $k$-colourable.
\emph{Independent $\Rightarrow q\notin\Acc(B)$.} If $|B|=1$, $\Acc(\{d_u\})=\{d_u\}\not\ni q$. If $|B|\ge2$, let $w$ be the least
vertex of $B$; all $d_u\in B$ have $S$ above depth $w$ on both sides, and at depth $w$ the datum $d_w$ has root $+$ and the
others $S$, so $\sigma_w,r_w$ are slots. For $u\in B$, $d_u|_{r_w}=(d_u|_{\sigma_w})[c_u]$. As no $y_i$ is free in two members
(independence), the restrictions of the $c_u$ are compatible, and their union $\bar u$ makes $\mathrm{Eq}(\sigma_w,r_w,\bar u)$
a $B$-feature. But $q|_{\sigma_w}=S^{n-w}(0)$ and $q|_{r_w}=S^{n+1-w}(0)$ are closed and different.
\emph{An edge $\{u,v\}\subseteq B\Rightarrow q\in\Acc(B)$.} $\Acc$ is monotone in the data, so it suffices that $q\in
\Feat(\{d_u,d_v\})$. With $w=\min(u,v)$, $C(\{d_u,d_v\})=\forall^m(S^w(\square)=0\wedge S^w(\square)=0)$ with slots
$\sigma_w,r_w$; $q$ has every Sym feature ($n>w$) and every Scope feature (its contents there are closed). Eq-features with
source $\sigma_w$: for $(\sigma_w,r_w)$, the edge's index $y_e$ is free at $\sigma_w$ in both data and $c_u(y_e)\neq c_v(y_e)$, so
by forcing there is no common $\bar u$; for the other admissible term positions $p\perp\sigma_w$ (the right-hand $0$'s and the
$S$-positions of the second conjunct above $r_w$), the datum with vertex $w$ has root $+$ at $\sigma_w$, and after any
substitution, but $0$ or $S$ at $p$. Eq-features with source $r_w$ have empty scope (the $c_u$ are closed), so they are
equalities; the datum with vertex $w$ has root $+$ at $r_w$, the candidates carry $0$ or $S$, except $\sigma_w$, which contains
a free index while $r_w$ is closed. So the only features are Sym and Scope features, which $q$ has.
\end{proof}

\begin{proof}[Proof of \cref{thm:many:complexity}(c)]
(i) Reduction from Set Cover \citep{karp1972reducibility}: universe $\{1,\dots,n\}$, sets $F_1,\dots,F_t$, bound $k$. Let
$d_u:=\bigwedge_{j\le t}(c_{uj}=0)$ with $c_{uj}=0$ if $u\in F_j$ and $S0$ otherwise, and $q:=\bigwedge_{j\le t}(S0=0)$. If some $u$
lies in no $F_j$, then $d_u=q\in D$, there is no cover, and $q\in\Acc_k(D)$, consistently. Otherwise read off the types of $q$
(\cref{lem:many:types}); let $\ell_j$ be the left-hand side of conjunct $j$, where data have $0$ or $S$ and $q$ has $S$. The
type $\mathrm{Sym}(\ell_j,0)$ has allowed set exactly $\{d_u:u\in F_j\}$. There are no other Sym-types (below an $S$ at
$\ell_j$ data and $q$ both have $0$) and no Scope-types (no indices). All matchers are empty maps, so an Eq-type needs
$q|_\sigma\neq q|_r$; the pairs with $q|_\sigma\neq q|_r$ that data can satisfy are $(\ell_j,\text{a }0\text{-position})$, its
converse, and $(\ell_j\cdot1,\ell_{j'})$ with $u\notin F_j$ (and its converse); each forces $d|_{\ell_{j'}}=0$ for some $j'$, so its
allowed set lies inside some $\{d_u:u\in F_{j'}\}$. Formula-sort pairs are void. So a block fails iff it lies inside some
$\{d_u:u\in F_j\}$, and $q\notin\Acc_k(D)$ iff $k$ of the sets cover the universe. All data are quantifier-free, so all covering
templates are first-order patterns.
(ii) For binder-free $D$ every matcher is the empty map, so Eq-types have no incompatible pairs and need no witness unless
$e_q$ exists, in which case it is compatible with every $e_d$ and the type is void. So every type has the form
``$B\subseteq$ allowed set'', there are polynomially many, and $q\notin\Acc_k(D)$ iff $D$ is covered by at most $k$ allowed sets,
which for fixed $k$ is decided by trying all $k$-tuples.
\end{proof}

\begin{proof}[Proof of \cref{thm:many:complexity}(d)]
\emph{In NP.} Some $T\in\Min(X)$ avoids $R$ iff some covering template does (it lies above a minimal one), iff some saturated
covering template does (minimal templates are saturated, \cref{thm:single:min}(c)). A saturated template is determined by an
occurrence set inside $A(X)$ and a pattern slot per metavariable, a certificate of polynomial size, checked by matching.
\emph{NP-hard} (from CNF-SAT; \citealp{garey1979computers}). Take variables $x_1,\dots,x_n$, distinct parameters
$a_1,\dots,a_n$, and $X=\{d_1,d_2\}$ with $d_1:=\forall x\bigwedge_j(x=a_j)$, $d_2:=\forall x\bigwedge_j(a_j=a_j)$. The common prefix
is everything except the left-hand sides $\sigma_j$ (column $(x,a_j)$, scope $\{x\}$); the right-hand sides $z_j$ carry the
rigid $a_j$. Besides Sym and Scope features, the $X$-features are exactly $\mathrm{Eq}(\sigma_j,z_j,x\mapsto a_j)$ ($\sigma_i\to
\sigma_j$ fails on $d_2$ as $a_i\neq a_j$; $\sigma_i\to z_j$ would need $a_i=a_j$; there are no formula slots). So the saturated
covering templates are $T_A:=\forall x\bigwedge_j(f_j(x)=\zeta_j)$, $A\subseteq[n]$, with $\zeta_j=f_j(a_j)$ for $j\in A$ and
$\zeta_j=a_j$ otherwise; they are pairwise incomparable ($f_j:=\lambda z.Sz$ separates them), so $\Min(X)=\{T_A\}$. For a clause $C$
(w.l.o.g.\ not containing both $x_j$ and $\neg x_j$) put $n_C:=\forall x\bigwedge_j\gamma_j$ with $\gamma_j=(x=a_j)$ if $x_j$ does not
occur in $C$, $\gamma_j=(Sx=a_j)$ if $x_j\in C$, and $\gamma_j=(Sx=Sa_j)$ if $\neg x_j\in C$. By matching, where $\gamma_j=(x=a_j)$ the
matcher takes $f_j:=\lambda z.z$ and both choices of $\zeta_j$ fit; where the left side is $Sx$ it takes $f_j:=\lambda z.Sz$, and
then $\zeta_j=a_j$ fits only $Sx=a_j$ while $\zeta_j=Sa_j$ fits only $Sx=Sa_j$. So $n_C\in\inst(T_A)$ iff the assignment ``$x_j$ true
iff $j\in A$'' falsifies $C$, and some $T\in\Min(X)$ avoids $R:=\{n_C:C\in\varphi\}$ iff $\varphi$ is satisfiable. The size is
$O(n|\varphi|)$. \emph{Polynomial cases.} Whether $X$ has a least covering template is decidable in polynomial time, and then
it is computed (\cref{prop:single:lgg}), and $\Min(X)$ is that template; under \Rich{} an anchor gives one
(\cref{cor:single:anchor-lgg}). If
$r\in R\cap\Acc(X)$, every covering template has $r$ as an instance.
\end{proof}
In this reduction $d_1$ is false, so the hardness does not transfer to \DTRC's use (true data, false negatives in a fixed
world); that case is open. Computed\src{single e11, e14}: the 2-SAT verifier equals the brute-force partition verifier on all
5{,}918 queries; for 12 graphs (including $K_3$, $C_5$, $K_4$, the Petersen graph) $q_G\in\Acc_k(D_G)$ iff $\chi(G)>k$ for
$k=1,\dots,4$; on 25 random set systems $q\in\Acc_k(D)$ iff no $k$ sets cover; for $n\le6$, $|\Min(X)|=2^n$ in the SAT
reduction, and the test equals satisfiability on 360/360 random CNFs. These results were not refereed independently
(\cref{app:ver:unref}).

%=====================================================================================================================
\subsection{MDL}\label{app:many:mdl}

MDL with likelihood errs in two directions. For schemas it errs towards incompleteness: a split is a sound but incomplete
hypothesis, which never accepts induction on an unused connective. For rarely used ground axioms it errs towards
unsoundness: two ground axioms used $m$ times each are merged by the naive code into their lgg once $m$ times the body
cost is below the template saving, and for $0+0=0$ and $0\cdot0=0$ that lgg, $z=0$, is false. The refutation test of
\DTRC{} catches the second error, and the anchor condition (data diversity) controls the first. This is the size principle
\citep{tenenbaum2001generalization} and MDL \citep{rissanen1978modeling} meeting a hypothesis space whose boundaries are
logical rather than statistical. The definitions and results follow.

For a finite set $H$ of templates and data $D$, each datum $d$ is assigned a covering template $T_d\in H$ (the cheapest if
several), with unique matcher $\theta_d$, and $L(H,D)=\sum_{T\in H}L_{\mathrm{tmpl}}(T)+\sum_{d\in D}[L_{\mathrm{idx}}(T_d)+
L_{\mathrm{body}}(\theta_d\mid T_d)]$. Without the body term this is the parent report's observation that MDL without
likelihood picks the bare metavariable \citep{claude2026whatfollows}. The split hypothesis is $H_F=\{Q1,\dots,Q7\}\cup\{T_f:f\in
F\}$ with $F$ the main connectives of the induction data and $T_f:=T_{\Ind}[P:=\lambda x.f(P_1(x),\dots)]$ (for ``$=$'':
$\lambda x.t_L(x)=t_R(x)$; for $\forall,\exists$: $\lambda x.\forall y\,P_1(x,y)$). For a datum $\Ind(\varphi)$ with root $f$ the matcher of
$T_f$ has bodies of total size $|\varphi|-1$; under $T_{\Ind}$ the body is $\lambda x.\varphi$.

\begin{proposition}[a naive code over-splits]\label{prop:many:mdlnaive}
\status{proved}\src{untagged Prop 6.2}
Let $L_{\mathrm{tmpl}}(T)=\lambda|T|$ and $L_{\mathrm{body}}(\theta)=\lambda\sum_M|\theta(M)|$ with $\lambda=\log_2A$, $A$ the number
of symbols; let $L_{\mathrm{idx}}$ be the two-part maximum-likelihood code of the label sequence (empirical entropy plus
$\frac{|H|-1}2\log_2n+O(1)$, e.g.\ the KT code \citep{krichevsky1981performance}); let $n_{\Ind}$ be the number of induction
data and $\hat H$ the empirical entropy of their main connectives. Then
\[
L(H_F,D)-L(H_{\mathrm{true}},D)=\lambda\Bigl(\sum_{f\in F}|T_f|-|T_{\Ind}|\Bigr)+\frac{|F|-1}2\log_2n-n_{\Ind}(\lambda-\hat H)+O(1).
\]
Since $\hat H\le\log_2|F|<\lambda$, the difference tends to $-\infty$ linearly in $n_{\Ind}$. With a uniform index code
($\log_2|H|$ bits per datum) the same holds once the induction frequency exceeds $\log_2((7+|F|)/8)/\lambda$. Applied
recursively, naive MDL keeps splitting each class by its next symbol while the class has enough data.
\end{proposition}
\begin{proof}
Template part: as displayed. Body part: each induction datum costs $\lambda$ fewer bits under its $T_f$, for every root,
$\forall$ and $\exists$ included. Index part: the split labels refine the true labels by the root of the induction data, so
their empirical entropy is $n\hat H(\text{true labels})+n_{\Ind}\hat H$ (chain rule), and the parametric term grows by
$\frac{|F|-1}2\log_2n$. Ground axioms cost the same under both hypotheses. For the uniform code, the index cost grows by
$n\log_2((7+|F|)/8)$ and the body cost falls by $n_{\Ind}\lambda$.
\end{proof}

\begin{proposition}[well-specified codes gain at most $O(\log n)$]\label{prop:many:mdlwell}
\status{proved}\src{untagged Prop 6.3}
Let the data be i.i.d.\ from a law $P_0$ (template index and instance law), and let the code for $H_{\mathrm{true}}$ be a
Bayesian mixture over an $r$-parameter family containing $P_0$ with regret at most $\frac r2\log_2n+C$
\citep{rissanen1986stochastic,clarke1990information} (a standard fact, cited from memory). Then for every fixed split hypothesis and $K>0$, for a constant
$C'$,
\[P\bigl[L(H_F,D)\le L(H_{\mathrm{true}},D)-\tfrac r2\log_2n-K-C'\bigr]\le2^{-K}.\]
\end{proposition}
\begin{proof}
The data part of $L(H_F,\cdot)$ is a prefix code for the data sequence given the template list, i.e.\ a sub-probability
$Q$; by the no-hypercompression inequality (see \citealp{grunwald2007minimum}), $P[-\log_2Q(D)\le-\log_2P_0(D)-K]\le2^{-K}$.
The data part of $L(H_{\mathrm{true}},\cdot)$ is at most $-\log_2P_0(D)+\frac r2\log_2n+C$, and the template parts are
constants; take $C'$ to absorb $C$ and the template parts.
\end{proof}
The regret bound and the no-hypercompression inequality are standard, but we cite them from memory and did not re-check
them. With a well-specified code, splitting cannot win linearly; the decision falls to lower-order terms, which favour
fewer templates. This does not prove that MDL selects the target partition. Computed (\cref{tab:many:mdl}): the naive code
splits from $n\approx10^3$; under the natural law G1 even the depth-aware code DPC splits for $n\ge1.6\cdot10^4$, because
usage really differs across connectives (terms under a quantifier see a bound variable); under G2, for which DPC is well
specified, DPC keeps $T_{\Ind}$ whole with a margin growing like $\log n$, while the misspecified PC eventually splits.

\begin{table}[ht]
\centering
\caption{$L(H_{\mathrm{root}})-L(H_{\mathrm{true}})$ in bits on PA data ($\Q$ and induction, induction frequency $0.5$);
negative means that MDL splits induction by main connective. NAIVE: \cref{prop:many:mdlnaive} with $A=23$; PC: sequential
KT per template with context (parent symbol, child index); DPC: KT with context (depth below the motive root, parent, child
index). G1: a natural usage law (random motives with quantifiers; scope-dependent terms); G2: a law for which DPC is well
specified (symbols depend only on depth, parent and index). Computed\src{untagged u7}, reproduced exactly by the
referee.}\label{tab:many:mdl}
\begin{tabular}{llrrr}
\toprule
usage law & $n$ & NAIVE & PC & DPC\\
\midrule
G1 natural & 1000 & $-223$ & $+2325$ & $+2701$\\
 & 16000 & $-16466$ & $-1199$ & $-3036$\\
 & 64000 & $-68243$ & $-19783$ & $-34412$\\
G2 well specified for DPC & 1000 & $-915$ & $+1141$ & $+1428$\\
 & 16000 & $-21467$ & $+1053$ & $+2822$\\
 & 64000 & $-87156$ & $-1522$ & $+3564$\\
\bottomrule
\end{tabular}
\end{table}

%=====================================================================================================================
\subsection{Natural examples: details}\label{app:many:examples}

\begin{proposition}[a passing run is \DTRC{} with a fixed oracle]\label{prop:many:audit}
\status{proved}\src{untagged Prop 5.12}
The implementation keeps a global set $N$ of refuted sentences; a template is reported refuted iff it covers some
$n\in N$ or its own budgeted search finds a refuted instance, which is then added to $N$. Every coherence decision is
logged. A run on data $D$ (fixed order and policy) ending with $N_{\mathrm{final}}$ \emph{passes the audit} if every
``coherent'' decision had a minimal template covering no element of $N_{\mathrm{final}}$, and every asserted template avoids
$N_{\mathrm{final}}$. If the sentence refuters are sound, then $\Ref_d:=N_{\mathrm{final}}$ satisfies (OS) and (Dec), and the
run's final clusters and asserted templates coincide with those of $\DTRC_d$ with this fixed oracle on the same data,
order and policy.
\end{proposition}
\begin{proof}
(OS) holds because $N_{\mathrm{final}}$ consists of sentences refuted by sound procedures; (Dec) holds by matching. The
merge policy is a deterministic function of the sequence of coherence answers, so it suffices that each answer agrees with
the fixed oracle: ``incoherent'' answers agree because $N\subseteq N_{\mathrm{final}}$ at every time (each refuted minimal
template covers its refuting sentence), and ``coherent'' answers agree by the audit. The asserted templates are the members
of $\Min(C)$ not reported refuted: every excluded member covers an element of $N_{\mathrm{final}}$, and every included one
avoids $N_{\mathrm{final}}$ by the audit.
\end{proof}
Refuter soundness is supported by code reading and by the referee's fuzzing (0 of 3000 true arithmetic and 1206 true
set-theoretic sentences refuted). Every \DTRC{} run passed the audit on the first pass (24 runs checked), and the re-audit
after the separation check found no inconsistent decision. (The prototype's first version decided refutation by
template-specific searches, which is not monotone in $\gen$; the global refuted set and the audit replaced it, and no
result changed, \cref{app:ver:corr}.) Caveats: completeness of the computed minimal
templates was checked against brute-force enumeration (by the author on induction pairs; by the referee's independent
enumerator on all PA and ZF cross pairs, several within-schema sets, the $K_0$ pair and the $J^*_0$ clusters), always with no
covering template outside the computed ones; internal caps never changed a result; templates are matched literally, not
modulo renaming of their parameters, which can only under-accept.

\begin{table}[ht]
\centering
\caption{PA cross merges: the unique minimal covering template of each cross pair and an R-$\Delta_0$ refutation (computed\src{untagged u1};
the referee's independent enumeration of all 68 covering templates of size $\le14$ of the 49 tested cross pairs,
21 $\Q$--$\Q$ and $7\times4$ $\Q$--induction, agrees). Numbering as in \cref{sec:setting:examples}, in parameter form. Q$_i^s$ is the instance
schema of Q$_i$ (closed terms for the parameters), an alternative presentation of $\Q$.}\label{tab:many:pacross}
\begin{tabular}{lll}
\toprule
cross pair & minimal template & refuting instance\\
\midrule
different roots (Q1--Q4 etc.); Q$_i$--Ind, root $\neq\to$ & $F_0$ & $0=S0$\\
Q2--Q3, Q2--Ind, Q3--Ind & $F_0\to F_1$ & $0=0\to0=S0$\\
Q4--Q6, Q4--Q7, Q5--Q6, Q5--Q7 & $z_0=z_1$ & $0=S0$\\
Q4--Q5 & $p+z_0=z_1$ & $p+0=0$ at $p:=1$\\
Q6--Q7 & $p\cdot z_0=z_1$ & $p\cdot0=S0$ at $p:=0$\\
Q$_i^s$--Q$_j^s$ (21 pairs, 4 samples each) & various & all refuted\\
\bottomrule
\end{tabular}
\end{table}

In the instance-schema presentation, one $\mathrm{Q2}^s$ sample lacked an $S$-free instance and gave the sound
specialization $SSz_0=Sz_1\to Sz_0=z_1$. The PA near-miss set had 104 sentences (50 step-by-two and 50 wrong-base induction
variants, 4 mutated axioms); 13 variants are genuine inductions (vacuous motives) and were accepted, the other 91 rejected.
Held-out data come from other seeds; the referee found no leakage. Separation on the data was verified for all 147, 201
and 175 cross pairs of the three PA runs.

\begin{example}[ZF cross merges]\label{ex:many:zfcross}
\status{computed}\src{untagged u4, u13; confirmed by the referee}
The ZF targets in closure-normal form: Ext $\forall x(x\in p\leftrightarrow x\in p')\to p=p'$; Pair $\exists c(p\in c\wedge p'\in c)$; Union
$\exists b\forall c(c\in b\leftrightarrow\exists d(d\in p\wedge c\in d))$; Power $\exists b\forall c(c\in b\leftrightarrow\forall d(d\in c\to d\in p))$;
Inf (empty set and successor written out with $\in$); Found $\exists y(y\in p)\to\exists y(y\in p\wedge\forall z(z\in y\to\neg z\in p))$;
Separation $\exists b\forall c(c\in b\leftrightarrow c\in p\wedge P(c))$; Replacement $\forall x(x\in p\to\exists y(P(x,y)\wedge\forall
z(P(x,z)\to z=y)))\to\exists b\forall y(y\in b\leftrightarrow\exists x(x\in p\wedge P(x,y)))$; $\in$-induction $\forall x(\forall y(y\in x\to
P(y))\to P(x))\to\forall xP(x)$. All 36 cross pairs (one instance per schema) have a unique minimal covering template:
$F_0$ or $F_0\to F_1$ in 22 pairs, refuted by an HF counterexample ($\neg p=p$, resp.\ $p=p\to\neg p=p$, at $p:=\emptyset$);
in 14 pairs a template keeping quantifier structure, refuted by logic: $\exists xF_0(x)$ (Pair--Union, Pair--Power,
Pair--Sep, Union--Inf, Power--Inf, Inf--Sep; instance $\exists x\neg p=p$), $\exists x(F_0(x)\wedge F_1(x))$ (Pair--Inf),
$(\forall xF_0(x))\to F_1$ (Ext--Rep, Ext--EInd), $F_0\to\exists xF_1(x)$ (Found--Rep), $(\forall x(F_0(x)\to F_1(x)))\to F_2$
(Rep--EInd), and naive comprehension $\exists x\forall y(y\in x\leftrightarrow F_0(y))$ (Union--Power, Union--Sep, Power--Sep), refuted
through Russell's instance (Herbrand: $y:=x$). The referee's independent enumerator listed all 598 covering templates of
size $\le14$ of 63 ZF cross pairs (two samples per schema): none outside the computed minimal ones, all refuted with the
global-negative-set oracle (the prototype's first version had left 120 of them unrefuted, all above refuted naive
comprehension; \cref{app:ver:corr}). Within each schema, three random instances have the schema as unique minimal template (confirmed by the
referee for Separation, $\in$-induction and Replacement).
\end{example}

\begin{proof}[Proof of \cref{prop:many:univset}]
Every instance holds in $M_0$ with $x:=u$; $\HF$ is a transitive part of $M_0$, so R-HF's steps ($\Delta_0$ facts about
$\HF$, HF witnesses and counterexamples) are sound in $M_0$; R-logic refutes only logical falsehoods. A witness $b$ of
$\exists x\forall y(y=y\to y\in x)$ gives $b\in b$; or $r=\{y\in b:y\notin y\}$ gives $r\in r\leftrightarrow r\notin r$.
\end{proof}

\begin{example}[\DTRC{} on ZF and ZFC]\label{ex:many:zf}
\status{computed}\src{untagged u5, u10}
\emph{ZF} (three seeds of 45 data): besides the numbers in \cref{sec:many:examples}, 60/60 held-out instances of each
schema were accepted, separation on the data was verified on all 467, 371 and 348 cross pairs, and $|\Min|=1$ in all 442
tests. \emph{ZFC}: in closure-normal form, with $p$ the family,
\begin{align*}
\AC(p)={}&[\forall y(y\in p\to\exists z\,z\in y)\wedge\forall y\forall w((y\in p\wedge w\in p\wedge\neg y=w)\to\neg\exists z(z\in y\wedge z\in w))]\\
&\to\exists c\forall y(y\in p\to\exists z(z\in y\wedge z\in c\wedge\forall u((u\in y\wedge u\in c)\to u=z))).
\end{align*} AC is true in $V$ and not
refuted. Its 12 cross pairs (two samples per schema) each have one minimal template: $F_0\to F_1$ (with Ext and EInd) and
$F_0$ (with Pair, Union, Power, Inf, Sep), refuted by an HF counterexample; $F_0\to\exists xF_1(x)$ (with Found) and
$F_0\to\exists x\forall yF_1(y,x)$ (with Rep), refuted by logic. On ZFC data (three seeds of 50): 10 pure clusters, schema clusters exact, AC
accepted, 180/180 held-out schema instances, separation verified on 614, 294 and 428 cross pairs. \emph{Weak Union and Power} (part B): with logic and HF only, the two weak
forms merge into the universal-set cluster and $\exists x\forall y(y=y\to y\in x)$ is accepted (8 clusters); with the closure of
Russell's Separation instance designated (bootstrapped coherence), 9 pure clusters and the universal set rejected. Over the
UnionW--PowerW pair, the referee's enumerator found 81 of 91 covering templates unrefuted, all above the universal-set
schema, as \cref{prop:many:univset} predicts.
\end{example}

\paragraph{$K_0$.} $K_0=(z_1+0=z_1)\wedge\forall x(z_1+x=z_2\to z_1+Sx=Sz_2)\to\forall x(z_1+x=z_2)$ is the first-order lgg
(named encoding) of $\Ind(0+x=x)$ and $\Ind(S0+x=Sx)$. By \cref{thm:zf:K0}, no sound refutation from quantifier-free truths
exists, so R-$\Delta_0$ and logic refute none of its instances, among them the false $J^*_0$. In $\DT$ the pair has two
minimal covering templates (computed; the referee's enumerator found no other among 6258 covering templates of size
$\le24$): $T_{\Ind}[P:=\lambda x.(f(0)+x=f(x))]$ and the $K_0$-shaped template of \cref{ex:many:noise}. The intersection
verifier excludes $J^*_0$ (it is not an instance of the first). The $K_0$-shaped template has $J^*_0$ as an instance
($f:=\lambda h.0$) and is refuted by coherence with designated $\Q$: the refutation of $J^*_0$ through the recursion axiom
Q5, $\forall x\forall y\,(x+Sy=S(x+y))$, goes through once Q5 is designated \src{prior induction B5(d)}
(see the text after \cref{thm:zf:K0}).

\paragraph{Computations.} The scripts are in \url{research/tracks/untagged/code} (Python 3.11, no external packages;
\url{sh run_all.sh}, about 15 minutes), with outputs \url{<script>.out}: u0 (minimal templates vs brute force), u1
(\cref{tab:many:pacross}), u2 ($c=4/7/7$), u3 (\DTRC{} on PA), u4, u5, u5b (ZF cross merges, \DTRC{} on ZF, weak forms), u6
($\forall x\varphi$ from numerals; non-transitivity), u7 (\cref{tab:many:mdl}), u8 (\cref{ex:many:noise}), u9 (clean
fragmentation), u10 (ZFC), u12 (head classes), u13, u14 (monotone oracle, audit). The checks of \cref{thm:many:complexity} are
\url{research/tracks/single/e11_union_complexity.py} and \url{e14_dtrc_step2.py}; the referee's independent code is in
\url{research/tracks/untagged/referee_code}.

\paragraph{Credits.} Gold's theorem \citep{gold1967language}; lgg and witness events
\citep{plotkin1970note,reynolds1970transformational}; finite elasticity of bounded unions
\citep{wright1989identification,motoki1991correct}; minimal multiple generalizations for unions of pattern languages
\citep{arimura1994finding} (overlap with \cref{thm:many:complexity} not checked); MRDP
\citep{davis1961decision,matiyasevich1993hilbert}; $\Delta_0$ absoluteness \citep{kunen1980set,jech2003set}; MDL and coding
\citep{rissanen1978modeling,rissanen1986stochastic,clarke1990information,krichevsky1981performance,grunwald2007minimum}; the size
principle \citep{tenenbaum2001generalization}; conceptual clustering \citep{michalski1983learning}, in whose spirit the
clustering is. Covering problems for first-order patterns, relevant to the open spare-slot thresholds, are decidable by
results on complement problems and ground reducibility \citep{lassez1987explicit,comon2003ground} (cited from memory, not
re-checked). As far as we know, the refutation test, the separation theorem, the depth impossibility, the head-class
thresholds and the implementation audit are new.
